Proposition \ref{prop.idealisedUpdate} is an identity about a window that has elapsed.
This subsection asks the forward question:
given everything observed up to $t$, where does the flow take the mid-price over
$(t, t+\forecastHorizon]$.
The flow half of the answer is exact, and it is the one we state first.

Throughout, the specification is direction-symmetric,
Assumption \ref{assumption.idealisedBook} holds over the window,
and $\sigmaAlgebra_t$ is the history of the order flow up to and including $t$.
We write $\decayedCounts(t+)$ for the state read after the event at $t$, if there is one.

\begin{prop}\label{prop.shortHorizonForecast}
 For every $\forecastHorizon \geq 0$,
 \begin{equation}\label{eq.shortHorizonForecast}
  \Expectation \big[ \OFI_{t, \forecastHorizon} \,\big|\, \sigmaAlgebra_t \big]
  = \meanSize \; \pressure^{\transpose} \excitation \,
  \meanResponseIntegral(\forecastHorizon) \, \decayedCounts(t+) ,
 \end{equation}
 where $\OFI_{t,\forecastHorizon}$ is the imbalance accumulated over
 $(t, t+\forecastHorizon]$ and
 $\meanResponseIntegral$ is as in Proposition \ref{prop.forwardMean}.
\end{prop}
\begin{proof}
 By Proposition \ref{prop.signedSizeContribution} the imbalance over the horizon is
 $\sum \pressure_{\event[n]}\orderSize_n$ over the events in $(t, t+\forecastHorizon]$.
 The marks are independent of the arrival times and of the types, with mean $\meanSize$,
 so conditioning on those,
 $\Expectation[\OFI_{t,\forecastHorizon} \mid \sigmaAlgebra_t]
 = \meanSize \, \pressure^{\transpose}
 \Expectation[\multiCountingProc(t+\forecastHorizon) - \multiCountingProc(t)
 \mid \sigmaAlgebra_t]$,
 the sum being absolutely summable because the expected count is finite.
 By Theorem \ref{thm.markovState} and Proposition \ref{prop.forwardMean}, that expectation is
 \begin{equation*}
  \forecastHorizon \, \baseIntensity
  + \excitation \, \meanResponseIntegral(\forecastHorizon) \, \decayedCounts(t+)
  + \excitation \Big( \int_{0}^{\forecastHorizon} \meanResponseIntegral(s) \, ds \Big)
  \baseIntensity .
 \end{equation*}
 The first and third terms are $\swapMatrix$-invariant matrices applied to $\baseIntensity$,
 so $\pressure^{\transpose}$ annihilates them by Proposition \ref{prop.directionSymmetry},
 and \eqref{eq.shortHorizonForecast} remains.
\end{proof}

\begin{corol}\label{corol.priceForecast}
 Under the same hypotheses,
 \begin{equation}\label{eq.priceForecast}
  \Expectation \big[ \midPrice_{t+\forecastHorizon} - \midPrice_{t}
  \,\big|\, \sigmaAlgebra_t \big]
  = \frac{\tickSizeOfLOB \meanSize}{2\uniformDepth} \;
  \pressure^{\transpose} \excitation \, \meanResponseIntegral(\forecastHorizon) \,
  \decayedCounts(t+)
  + \Expectation \big[ \residual_{t,\forecastHorizon} \,\big|\, \sigmaAlgebra_t \big] ,
 \end{equation}
 with $\vert \residual_{t,\forecastHorizon} \vert \leq \tickSizeOfLOB$.
\end{corol}
\begin{proof}
 Take conditional expectations in \eqref{eq.idealisedUpdate} over
 $(t, t+\forecastHorizon]$ and apply \eqref{eq.shortHorizonForecast}.
\end{proof}

\begin{remark}\label{remark.residualAgainstForecast}
 The residual is carried rather than absorbed, for two reasons.
 By Remark \ref{remark.residualIsQueueFlow} it is not independent of the first term:
 over a window on which neither quote moves it cancels that term exactly.
 And it does not shrink as the forecast grows:
 $\meanResponseIntegral(\forecastHorizon) \to -\hurwitzMatrix^{-1}$ as
 $\forecastHorizon \to \infty$ when $\branchingRatio < 1$,
 so the first term of \eqref{eq.priceForecast} is bounded in $\forecastHorizon$,
 while the residual is bounded by one tick at every horizon.
 A long horizon therefore does not buy a forecast that exceeds the residual.
 What does is a large state:
 \eqref{eq.priceForecast} carries information about the mid-price only where
 $\meanSize \,\vert \pressure^{\transpose}\branchingMatrix (I - \branchingMatrix)^{-1}
 \decayedCounts(t+) \vert$ is itself of the order of $\uniformDepth$.
 Section \ref{sec.lobsterEmpirics} asks the corresponding empirical question about the sign
 of \eqref{eq.priceForecast}, and not about its magnitude.
\end{remark}

\begin{corol}\label{corol.contrastAsDrift}
 As $\forecastHorizon \to 0$,
 \begin{equation}\label{eq.contrastAsDrift}
  \Expectation \big[ \OFI_{t,\forecastHorizon} \,\big|\, \sigmaAlgebra_t \big]
  = \meanSize \, \intensityContrast(t+) \, \forecastHorizon
  + O(\forecastHorizon^{2}) .
 \end{equation}
\end{corol}
\begin{proof}
 $\meanResponseIntegral(\forecastHorizon)
 = \forecastHorizon I + \tfrac{1}{2}\forecastHorizon^{2}\hurwitzMatrix
 + O(\forecastHorizon^{3})$,
 and $\pressure^{\transpose}\excitation\decayedCounts(t+) = \intensityContrast(t+)$ by
 \eqref{eq.contrastFromState}.
\end{proof}

The intensity contrast is therefore the drift of the order flow imbalance.
That is what the contrast is for, and it is a statement to leading order in the horizon and
about the flow, the price inheriting it only through \eqref{eq.priceForecast}.

\begin{remark}\label{remark.totalResponse}
 At the other end, if $\branchingRatio < 1$ then
 $\excitation \meanResponseIntegral(\infty)
 = -\excitation\hurwitzMatrix^{-1}
 = \branchingMatrix (I - \branchingMatrix)^{-1}$,
 which is the whole expected progeny of the present state:
 direct offspring composed with all descendants.
 Hurwitz-ness is Proposition \ref{prop.hurwitzCriterion}, and uses $\excitation \geq 0$.
 By Definition \ref{def.relaxationTime} the state-dependent part of
 \eqref{eq.shortHorizonForecast} is still moving at horizons short against
 $1/(\decay(1-\branchingRatio))$, and has finished at horizons long against it.
\end{remark}

\begin{corol}\label{corol.regimeSign}
 Let $\branchingRatio < 1$ and set
 \begin{equation}\label{eq.signedExcitationAtHorizon}
  \signedExcitation_{\eone}(\forecastHorizon)
  := \frac{\big( \pressure^{\transpose}\excitation
  \meanResponseIntegral(\forecastHorizon) \big)_{\eone}}{\pressure_{\eone}} .
 \end{equation}
 Then $\signedExcitation(\forecastHorizon)$ takes the same value on the two members of each
 pair, and for a state carrying one extra event of a type $a$ with $\pressure_a = +1$
 against one fewer of its mirror,
 the forecast \eqref{eq.shortHorizonForecast} moves by
 $2\meanSize\signedExcitation_{a}(\forecastHorizon)$.
 Moreover $\signedExcitation(\forecastHorizon)/\forecastHorizon \to \signedExcitation$ as
 $\forecastHorizon \to 0$.
\end{corol}
\begin{proof}
 $\excitation\meanResponseIntegral(\forecastHorizon)$ is $\swapMatrix$-invariant,
 so the argument of Proposition \ref{prop.signedExcitation} applies verbatim;
 the displacement is
 $\pressure^{\transpose}\excitation\meanResponseIntegral(\forecastHorizon)
 (\unitVector_a - \unitVector_{\swapMatrix(a)})
 = \pressure_a \signedExcitation_a - \pressure_{\swapMatrix(a)}
 \signedExcitation_{\swapMatrix(a)}$ at horizon $\forecastHorizon$,
 which is $2\signedExcitation_a(\forecastHorizon)$.
\end{proof}

The sign of $\signedExcitation(\forecastHorizon)$ therefore decides whether the flow is
forecast to continue in the direction it has been pushing or to reverse,
and that sign belongs to the kernel.
Two specifications with the same branching ratio and the same stationary intensity can carry
$\signedExcitation$ of opposite sign, and then they forecast in opposite directions at every
horizon at which the signs persist,
while Proposition \ref{prop.idealisedUpdate} holds identically in both.
No statement of the form ``the order flow imbalance predicts continuation'' is made in these
notes without naming the specification it is made in.

\begin{remark}\label{remark.imbalanceAsProxy}
 Three statistics of the same past flow have appeared, and they differ in two ways.
 The signed count $\pressure^{\transpose}
 \big( \multiCountingProc(t) - \multiCountingProc(t-\ofiWindow) \big)$
 weights the events of a window equally and ignores their sizes;
 the exponentially weighted count $\pressure^{\transpose}\decayedCounts(t)$
 weights them by age and ignores their sizes;
 and $\OFI_{t,\ofiWindow}$ weights them equally and by their sizes.
 The forecast \eqref{eq.shortHorizonForecast} is driven by none of the three,
 but by $\pressure^{\transpose}\excitation\decayedCounts(t+)$,
 which weights them by age and by $\signedExcitation$.

 The gap in the weighting closes exactly when $\signedExcitation$ is constant across the six
 types, that is when $\pressure$ is a left eigenvector of $\excitation$:
 then $\pressure^{\transpose}\excitation\meanResponseIntegral(\forecastHorizon)$ is a
 multiple of $\pressure^{\transpose}$ at every horizon,
 and the forecast is a multiple of $\pressure^{\transpose}\decayedCounts(t+)$.
 The gap in the marks does not close at all, by the last sentence of
 Section \ref{sec.readingTheIntensity}.
 So $\OFI$ is a statistic of the flow that needs no parameters and estimates the forecast
 only up to those two gaps, and how much is lost to each is a question about a particular
 specification.
\end{remark}
